% Kandydat restrukturyzacji; nie jest aktywnym rozdziałem Atlasu.

\section{Operator maksymalny Hardy'ego--Littlewooda}
\label{sec:dt-maximal}

Rozdział opiera się na
\atlasref{ME-24::thm:HL-weak-24}{ME-24: nierówności słabego typu dla funkcji maksymalnej}.
Używamy również lematu o pięciokrotnym powiększeniu
\atlasref{ME-24::thm:5r-24}{ME-24: lemat 5r}.
Dla $d\ge1$ niech $v_d=\lambda_d(B(0,1))$.

\begin{definicja}[Funkcja maksymalna Hardy'ego--Littlewooda]
\label{def:dt-HL-max}
Dla $f\in L^1_{\mathrm{loc}}(\mathbb R^d)$ definiujemy
\emph{centrowaną funkcję maksymalną}
\[
 Mf(x):=\sup_{r>0}\frac1{\lambda_d(B(x,r))}
 \int_{B(x,r)}|f(y)|\,\mathrm d\lambda_d(y)
 \in[0,+\infty].
\]
Definicja nie zależy od reprezentanta klasy $f$ w
$L^1_{\mathrm{loc}}$, ponieważ całki funkcji równych
prawie wszędzie są równe na każdej kuli.
Funkcja $Mf$ jest mierzalna jako supremum przeliczalnej
rodziny ciągłych średnich dla dodatnich promieni wymiernych. Dokładniej, gdy $(x_n,r_n)\to(x,r)$ z $r>0$,
wszystkie kule $B(x_n,r_n)$ od pewnego indeksu
mieszczą się w jednej ograniczonej kuli $D$.
Dla $y\notin\partial B(x,r)$ zachodzi
$\chi_{B(x_n,r_n)}(y)\to\chi_{B(x,r)}(y)$.
Sfera ma miarę zero, a funkcja $|f|\chi_D$
jest całkowalna. Twierdzenie Lebesgue'a o zbieżności
zmajoryzowanej daje ciągłość licznika; mianownik
$v_dr_n^d\to v_dr^d>0$. W szczególności można
ograniczyć supremum do dodatnich promieni wymiernych.
\end{definicja}

\begin{twierdzenie}[Nierówność słabego typu $(1,1)$]
\label{thm:dt-HL-weak}
Jeśli $f\in L^1(\mathbb R^d)$, to dla każdego $t>0$
\[
 \lambda_d\bigl(\{x\in\mathbb R^d:Mf(x)>t\}\bigr)
 \le\frac{5^d}{t}\int_{\mathbb R^d}|f(x)|\,
 \mathrm d\lambda_d(x).
\]
\end{twierdzenie}
\begin{proof}
Jeśli $\{Mf>t\}=\varnothing$, teza jest oczywista.
W przeciwnym razie ustalmy niepusty zwarty zbiór
$K\subseteq\{Mf>t\}$ (dla pustego $K$ żądane
oszacowanie jest natychmiastowe).
Dla każdego $x\in K$ wybieramy promień $r_x>0$ taki,
że $\int_{B(x,r_x)}|f|\,\mathrm d\lambda_d>
 t\lambda_d(B(x,r_x))$. Promienie takich kul są
ograniczone z góry: ponieważ ich objętość jest mniejsza
od $\|f\|_1/t$, zachodzi $r_x^d<\|f\|_1/(tv_d)$.
Z przytoczonego lematu $5r$ wybieramy rozłączne kule $(B_j)$, których
pięciokrotne powiększenia pokrywają $K$. Wówczas
\[
 \lambda_d(K)\le5^d\sum_j\lambda_d(B_j)
 <\frac{5^d}{t}\sum_j\int_{B_j}|f|\,\mathrm d\lambda_d
 \le\frac{5^d}{t}\|f\|_1.
\]
Zbiór $\{Mf>t\}$ jest otwarty: dla dowolnego $x$
warunek średniej większej od $t$ zachowuje się przy małej
zmianie środka kuli dzięki ciągłości całki z $|f|$.
Jego miara jest supremum miar zwartych podzbiorów,
na mocy regularności miary Lebesgue'a. Biorąc supremum
po $K$, kończymy dowód.
\end{proof}

\begin{wniosek}[Kontrola zbioru, na którym średnie są duże]
\label{cor:dt-HL-level}
Dla $f\in L^1(\mathbb R^d)$ oraz $a>0$ funkcja $Mf$
jest skończona prawie wszędzie, a
$\lambda_d(\{Mf>a\})\le5^d\|f\|_1/a$.
\end{wniosek}
\begin{proof}
Oszacowanie wynika z twierdzenia~\ref{thm:dt-HL-weak}.
Przechodząc do granicy $a\to\infty$, dostajemy
$\lambda_d(\{Mf=+\infty\})=0$.
\end{proof}


\subsection{Ograniczoność w przestrzeniach Lebesgue'a}

\begin{twierdzenie}[Silne oszacowanie operatora maksymalnego]
\label{thm:dt-HL-strong}
Niech $1<p<\infty$ i $f\in L^p(\R^d)$.
Wówczas $Mf\in L^p(\R^d)$ oraz
\[
 \|Mf\|_p^p\le \frac{2^p5^d p}{p-1}\|f\|_p^p.
\]
Dla $p=\infty$ zachodzi $\|Mf\|_\infty\le\|f\|_\infty$.
\end{twierdzenie}
\begin{proof}
Dla $1<p<\infty$ nierówność Höldera na każdej
ograniczonej kuli zapewnia $f\in L^1_{\mathrm{loc}}$.
Zatem $Mf$ jest określona i mierzalna na mocy
definicji \ref{def:dt-HL-max}. Dla $t>0$ rozkładamy
$f=f_1+f_2$, gdzie
$f_1=f\chi_{\{|f|>t/2\}}$ i $f_2=f-f_1$.
Z nierówności $|f_1|\le(2/t)^{p-1}|f|^p$
wynika $f_1\in L^1$. Ponadto $|f_2|\le t/2$
prawie wszędzie, więc średnia modułu $f_2$
na każdej kuli nie przekracza $t/2$.
Nierówność trójkąta pod całką i supremum po promieniach
dają $Mf\le Mf_1+Mf_2\le Mf_1+t/2$.
Stąd $\{Mf>t\}\subseteq\{Mf_1>t/2\}$.
Z twierdzenia \ref{thm:dt-HL-weak} otrzymujemy
\[
 \lambda_d(\{Mf>t\})
 \le \frac{2\,5^d}{t}
 \int_{\{|f|>t/2\}}|f(x)|\,\mathrm d\lambda_d(x).
\]
Dla każdej liczby $a\in[0,+\infty]$ zachodzi
$a^p=\int_0^\infty p t^{p-1}\chi_{\{t<a\}}\,\mathrm dt$;
dla $a=+\infty$ rozumiemy obie strony jako $+\infty$.
Podstawiając $a=Mf(x)$ i stosując twierdzenie Tonellego
do nieujemnej funkcji mierzalnej, otrzymujemy
\[
 \begin{aligned}
 \int_{\R^d}(Mf)^p\,\mathrm d\lambda_d
 &=p\int_0^\infty t^{p-1}
            \lambda_d(\{Mf>t\})\,\mathrm dt\\
 &\le 2\,5^d p\int_{\R^d}|f(x)|
       \left(\int_0^{2|f(x)|}t^{p-2}\,\mathrm dt\right)
       \mathrm d\lambda_d(x)\\
 &=\frac{2^p5^d p}{p-1}\int_{\R^d}|f|^p\,\mathrm d\lambda_d.
 \end{aligned}
\]
W drugiej linii ponownie stosujemy Tonellego;
całka wewnętrzna jest skończona dla skończonego
$|f(x)|$, ponieważ $p-2>-1$. Ostatnia strona
jest skończona, co daje tezę i skończoność $Mf$
prawie wszędzie.
Dla $p=\infty$ średnia na każdej kuli jest
ograniczona przez $\|f\|_\infty$, a supremum
zachowuje nierówność.
\end{proof}

\begin{uwaga}
Oszacowania maksymalne i twierdzenie Lebesgue'a
o różniczkowaniu całek
(\atlasref{ME-25::thm:lebesgue-diff-25}{ME-25: różniczkowanie całek})
są podstawą dalszego rozkładu Calderóna--Zygmunda.
Przypadek $p=1$ wymaga słabego oszacowania;
nie wynika z powyższej nierówności silnego typu.
\end{uwaga}
